% ---
% Capítulo 2: Fundamentação Teórica, Trabalhos Relacionados e Metodologia
% ---

\chapter{Fundamentação Teórica}

A evolução dos sistemas de informação nas últimas décadas foi acompanhada por transformações significativas nos padrões de arquitetura de \textit{software}. Com o crescimento da complexidade das aplicações e a exigência de entregas contínuas e escaláveis, os modelos tradicionais de desenvolvimento baseados em aplicações monolíticas passaram a apresentar limitações estruturais significativas. Nesse cenário, o estilo arquitetural de microsserviços consolidou-se como uma abordagem proeminente para a construção de sistemas modernos, modulares e desacoplados.

\section{Arquitetura Monolítica}

Historicamente, a maioria dos sistemas de \textit{software} foi desenvolvida sob o paradigma monolítico. De acordo com \citeonline{richardson2018}, uma aplicação monolítica é concebida como uma unidade lógica singular de execução. Em um monólito típico, todos os componentes do sistema, tais como a interface com o usuário, a lógica de negócio e os módulos de acesso a dados, são empacotados e implantados conjuntamente em um único processo computacional.

Embora a arquitetura monolítica apresente vantagens nos estágios iniciais do ciclo de vida do \textit{software}, como a simplicidade de testes locais, a facilidade de implantação inicial e a alta performance de chamadas de função diretas em memória, ela frequentemente enfrenta severos gargalos à medida que o sistema e a organização expandem-se, conforme cita \cite{newman2021}. 

Um dos principais problemas decorre do alto acoplamento de código e de banco de dados. A proliferação de dependências diretas entre módulos e o compartilhamento de um único esquema de banco de dados relacional tornam o sistema rígido. Nessa estrutura, pequenas alterações em uma funcionalidade específica podem desencadear efeitos colaterais imprevistos em módulos aparentemente não relacionados \cite{richardson2018}.

Adicionalmente, a arquitetura monolítica impõe uma limitação severa à escalabilidade granular. Para escalar uma aplicação monolítica sob alta demanda, é necessário duplicar a unidade inteira de execução em múltiplos servidores. Esse processo consome recursos de infraestrutura de forma ineficiente, visto que módulos ociosos são replicados juntamente com aqueles que realmente necessitam de maior capacidade de processamento \cite{lewis2014}.

Por fim, manifesta-se o acoplamento no ciclo de implantação. Toda a aplicação precisa ser recompilada, testada integralmente e reimplantada a cada nova atualização, por menor que seja a alteração no código-fonte. Esse fator eleva substancialmente a latência de entrega de novas funcionalidades e aumenta o risco operacional de falha global, onde uma exceção não tratada em um módulo secundário pode derrubar todo o processo da aplicação.

\section{Arquitetura de Microsserviços}

Para suplantar as limitações estruturais do monólito, \citeonline{lewis2014} definem a arquitetura de microsserviços como uma abordagem em que a aplicação é projetada como um ecossistema de pequenos serviços autônomos. Cada serviço é executado em seu próprio processo computacional e se comunica com os demais por meio de mecanismos leves de rede, tipicamente utilizando APIs fundamentadas no protocolo HTTP com o padrão REST ou barramentos de mensageria assíncrona.

Segundo as diretrizes estabelecidas por \citeonline{newman2021}, os microsserviços são unidades de \textit{software} independentemente implantáveis, fortemente coesas e estruturadas em torno de limites de negócios específicos. Em vez de organizar o código por camadas técnicas horizontais, tais como camada de apresentação, camada de regras de negócio e camada de persistência, os microsserviços alinham a estrutura do \textit{software} às capacidades de negócio da organização, promovendo o desacoplamento e a clareza de responsabilidade de cada componente.

\subsection{Características Fundamentais}

A literatura especializada destaca que a transição para microsserviços fundamenta-se em um conjunto de características e princípios arquiteturais estruturantes. A primeira dessas características refere-se à componentização promovida por meio de serviços. Conforme enfatizado por \citeonline{lewis2014}, a componentização tradicional via bibliotecas ligadas em memória é substituída por serviços independentes expostos via rede. Essa mudança assegura que a substituição, atualização ou manutenção de um serviço específico ocorra sem a necessidade de recompactar ou reimplantar os demais componentes do ecossistema.

Em adição à substituição de bibliotecas por serviços de rede, sobressai-se a governança e a descentralização tecnológica. Diferente do monólito, que impõe uma única linguagem e um único \textit{framework} para toda a aplicação, \citeonline{lewis2014} ressaltam que os microsserviços viabilizam a adoção da abordagem da "ferramenta certa para o trabalho", abrindo espaço para o que \citeonline{richardson2018} descreveu historicamente como uma "arquitetura poliglota". Sob essa ótica, \citeonline{newman2020} reforça que cada microsserviço é tecnologicamente agnóstico, garantindo autonomia para escolher a linguagem de programação, o \textit{framework} e as bibliotecas que melhor se adaptem às suas exigências técnicas particulares.

No que tange ao gerenciamento da informação, o isolamento de dados por serviço é apontado por \citeonline{richardson2018} como um pilar incontornável. Cada microsserviço deve possuir e gerenciar seu próprio repositório de dados privado, sendo impedido o acesso direto de serviços externos ao seu banco de dados. Qualquer interação ou consulta a informações sob domínio de outro serviço deve ocorrer exclusivamente por meio de sua API pública. Esse princípio elimina o acoplamento oculto no nível de esquemas relacionais de banco de dados.

Dessa segregação de dados deriva diretamente a possibilidade de adoção da persistência poliglota. Como o banco de dados deixa de ser um recurso compartilhado globalmente, \citeonline{richardson2018} explica que é possível selecionar tecnologias de armazenamento distintas de acordo com a natureza de cada domínio de negócio. Por exemplo, um serviço de catálogo de produtos pode adotar um banco NoSQL orientado a documentos para otimização de leitura, enquanto o serviço de carrinho de compras utiliza um banco chave-valor em memória para alta performance, e o serviço de pedidos utiliza um banco relacional garantidor de transações ACID \cite{richardson2018}.

Por fim, a separação física entre processos contribui de maneira determinante para a resiliência operacional e o isolamento de falhas. Em uma arquitetura monolítica tradicional, uma exceção não tratada pode paralisar todo o processo computacional. Em contrapartida, sob o modelo de microsserviços, a degradação ou a indisponibilidade pontual de um serviço específico não provoca o colapso de toda a aplicação. A adoção de padrões de resiliência impede a propagação de falhas em cascata, permitindo que funcionalidades não afetadas continuem operando normalmente para o usuário final \cite{newman2021}.

\subsection{Desafios da Arquitetura de Microsserviços}

Apesar dos evidentes benefícios de disponibilidade e desacoplamento, a adoção de microsserviços introduz uma nova camada de complexidade inerente aos sistemas distribuídos \cite{lewis2014}. Conforme adverte \citeonline{richardson2018}, a principal contrapartida operacional reside na observabilidade e no diagnóstico de falhas. Enquanto em um monólito a depuração ocorre pela análise de um log único e unificado com a pilha de execução completa (\textit{stack trace}), em microsserviços a execução de uma única requisição de usuário atravessa múltiplos serviços independentes via rede.

Essa pulverização de processos torna a identificação das causas raízes de erros um desafio complexo. Para contornar essa dificuldade e correlacionar os registros de auditoria espalhados entre diversos \textit{containers}, torna-se indispensável a implementação de técnicas de rastreamento distribuído (\textit{distributed tracing}), injetando identificadores únicos de correlação em todas as chamadas de rede \cite{richardson2018}. Além disso, destacam-se como desafios a latência de rede adicional, a exigência de roteamento eficiente por meio de \textit{API Gateways} e o gerenciamento de transações entre múltiplos repositórios.

Em particular, a substituição de transações ACID locais por modelos de consistência eventual exige a reformulação do fluxo de dados. A coordenação de estados entre serviços passa a demandar padrões avançados, tais como a Arquitetura Baseada em Eventos (\textit{Event-Driven Architecture}) e o Padrão Saga \cite{richardson2018}. Dessa forma, a transição para microsserviços representa uma escolha consciente (\textit{trade-off}): reduz-se a complexidade interna do código-fonte e ganha-se resiliência, em troca do aumento da complexidade de infraestrutura, rastreabilidade e comunicação de rede.

\section{Padrões de Decomposição}

Serão utilizados três diferentes padrões 


\subsection{Decomposição por Capacidade de Negócio}
% Visão macro e processual.

A decomposição por capacidade de negócio (\textit{Decompose by Business Capability}) fundamenta-se na identificação das funções operacionais e processos que a organização executa para atingir seus objetivos \cite{richardson2018}. Nesse contexto, uma capacidade de negócio representa o objetivo funcional da organização (o "que" é realizado), em contraste com as tecnologias usadas no momento (o "como" é realizado).

Dessa forma, a conceituação de uma capacidade permanece estável ao longo do tempo, independentemente das transformações nas ferramentas de \textit{software} ou na infraestrutura. Por exemplo, a capacidade de "Gestão de Estoque" ou de "Processamento de Pedidos" representa uma necessidade operacional contínua de uma organização comercial, quer seja executada por meio de registros manuais, por um sistema monolítico legado ou por microsserviços isolados.

Nessa estratégia, as fronteiras dos microsserviços são alinhadas às áreas funcionais da organização. Em um sistema corporativo, por exemplo, funções como gestão de estoque, processamento de pedidos e faturamento dão origem a serviços independentes. A principal vantagem teórica dessa abordagem consiste no alinhamento direto entre a estrutura do \textit{software} e a estrutura organizacional, o que favorece a atribuição de equipes autônomas responsáveis pelo ciclo de vida completo de cada serviço \cite{newman2021}.

Entretanto, essa abordagem apresenta limitações quando aplicada a sistemas monolíticos legados que utilizam um banco de dados relacional compartilhado. Como diferentes processos de negócio frequentemente acessam e alteram as mesmas tabelas, a divisão dos serviços baseada apenas na capacidade funcional pode exigir a implementação de transações distribuídas e gerar dependências cruzadas entre os novos serviços.



\subsection{Decomposição por Subdomínio (DDD)}
% Visão de fronteiras lógicas e isolamento de dados.

A decomposição por subdomínio baseia-se nos princípios do \textit{Domain-Driven Design} (DDD), propostos por \citeonline{evans2003} e aplicados à arquitetura de microsserviços por \citeonline{richardson2018}. Nessa abordagem, o domínio do sistema é particionado em subdomínios, classificados em centrais, de suporte e genéricos, sendo cada um associado a um Contexto Delimitado (\textit{Bounded Context}).

O conceito de \textit{Bounded Context} estabelece que cada modelo de dados e vocabulário técnico possui significado estrito dentro de seus limites lógicos \cite{evans2003}. Por exemplo, a entidade "Cliente" no contexto de vendas possui atributos e regras de negócio distintos da mesma entidade no contexto de atendimento ou faturamento.

No contexto experimental deste trabalho, a decomposição por subdomínio orienta a separação física dos serviços com base no isolamento dos modelos de dados. Essa estratégia busca reduzir o acoplamento conceitual e assegurar que cada microsserviço mantenha a consistência de seu modelo dentro de um repositório de dados privado, eliminando o acesso direto a bancos de dados compartilhados \cite{newman2020}.


\subsection{Decomposição por Transações}
% Visão de consistência técnica (ACID).

A decomposição por transações (\textit{Decompose by Transactions}) adota um critério técnico focado na consistência dos dados para a definição das fronteiras dos microsserviços \cite{richardson2018}. Diferente das abordagens orientadas a processos ou domínios de negócio, essa estratégia prioriza a análise das restrições ACID (\textit{Atomicity, Consistency, Isolation, Durability}) e do fluxo de operações de escrita no banco de dados.

Na refatoração de uma aplicação monolítica, a divisão de um banco de dados relacional centralizado pode transformar transações locais em transações distribuídas entre diferentes serviços. A coordenação dessas transações via rede adiciona complexidade e reduz a vazão do sistema, exigindo o uso de padrões como Saga ou commit de duas fases (\textit{Two-Phase Commit - 2PC}) \cite{richardson2018, newman2021}.

Para evitar a degradação de desempenho e a complexidade operacional, a decomposição por transações agrupa no mesmo microsserviço os módulos e tabelas que exigem escrita atômica e consistência imediata. Operações que suportam consistência eventual são isoladas em serviços assíncronos. Na matriz experimental desta pesquisa, essa abordagem atua como contraponto técnico voltado à redução de chamadas de rede e à preservação da integridade transacional.

\section{Métricas e Instrumentos de Coleta}
\label{sec:metricas_coleta}
% Apresentar as ferramentas (Docker stats, GitHub Actions, Autocannon, Madge/ESLint).

\subsection{Métricas de Agilidade e DevOps}


Para avaliar o impacto da decomposição na operabilidade e no ciclo de vida do software, esta dimensão foca na agilidade de entrega e na inicialização dos serviços:
\begin{itemize}
    \item \textbf{Tempo de Construção:} Mensura a duração total dos processos de empacotamento e geração das imagens Docker para cada microsserviço ou snapshot, avaliando a granularidade do pipeline de CI/CD.
    \item \textbf{Tempo de Inicialização:} Registra o intervalo necessário desde o comando de execução até a prontidão completa do serviço para receber requisições de rede (prontidão de escuta nas portas configuradas).
\end{itemize}

\subsection{Métricas de Custo Operacional e Infraestrutura}
\label{sec:metricas_infraestrutura}

O consumo de recursos computacionais é monitorado de forma isolada por contêiner para mensurar o \textit{overhead} introduzido pela distribuição da arquitetura:
\begin{itemize}
    \item \textbf{Uso de CPU e Memória:} Coletados por meio de ferramentas de monitoramento de contêineres, registrando a média e o pico de consumo de processamento e RAM durante a execução sob carga.
    \item \textbf{Uso de Armazenamento:} Avalia o impacto do particionamento e da duplicação de dependências analisando o tamanho final das imagens Docker geradas para cada componente da arquitetura.
\end{itemize}

\subsection{Métricas de Desempenho}
% Latência de API e Vazão Máxima.

Para avaliar o comportamento dinâmico e o impacto de gargalos operacionais gerados pela transição arquitetural, a pesquisa emprega métricas de desempenho, com ênfase na latência de API e na vazão máxima (\textit{throughput}). O fundamento empírico dessas medições baseia-se na identificação de restrições de desempenho provocadas pela conversão de dependências locais em chamadas de rede. Conforme apontam \citeonline{richardson2018} e \citeonline{newman2020}, a introdução da comunicação interprocessos (IPC) nas fronteiras dos microsserviços adiciona sobrecarga temporal que afeta diretamente a eficiência do sistema.

A latência é quantificada por meio de percentis, que revelam o comportamento sob estresse e a estabilidade nos tempos de resposta da aplicação. Ademais, a vazão máxima mensura o volume limite de requisições processadas por segundo sob condições controladas de concorrência. Nesse sentido, essas métricas permitem correlacionar diretamente a eficiência operacional da aplicação com o grau de acoplamento distribuído e a sobrecarga de rede imposta por cada abordagem de decomposição. 



\subsection{Métricas de Manutenibilidade e Complexidade}
\label{sec:metricas_manutenibilidade}

Para avaliar a evolução estrutural da arquitetura ao longo da refatoração, a pesquisa emprega métricas voltadas à manutenibilidade e complexidade, com destaque para as Linhas de Código (LOC), a contagem de componentes físicos e o cruzamento de fronteiras. O fundamento teórico para a análise desse acoplamento baseia-se no conceito de conascência (\textit{connascence}), definido formalmente como a interdependência na qual a alteração em um elemento de \textit{software} exige a modificação de outro para preservar a corretude do sistema \cite{pagejones1992}. O objetivo central do \textit{design} de arquiteturas modulares é a minimização da conascência e do cruzamento de fronteiras entre domínios distintos.

No monólito, o cruzamento de fronteiras é observado como conascências estáticas de nome e tipo, quantificadas por meio de análise estática do código-fonte. Durante a transição para microsserviços por meio do padrão \textit{Strangler Fig}, a remoção dessas dependências físicas evita o acoplamento indesejado que, conforme alertam \citeonline{newman2020} e \citeonline{richardson2018}, converte simples chamadas locais de código em comunicações síncronas via rede (\textit{Inter-Process Communication} - IPC). Além disso, o rastreamento das Linhas de Código (LOC) e do número de componentes físicos permite mensurar o esforço de refatoração, o dimensionamento dos artefatos e o impacto do isolamento em cada estágio da pesquisa.

\section{Ferramentas}
\subsection{Docker}
A plataforma de conteinerização \textit{Docker} foi adotada para garantir o isolamento consistente dos ambientes de execução em todas as fases da pesquisa. Além de prover portabilidade e padronizar o empacotamento dos serviços, o ambiente conteinerizado viabiliza a coleta controlada de métricas de infraestrutura, permitindo extrair o consumo computacional de CPU e o uso de memória RAM por meio do comando \textit{docker stats} durante os testes experimentais.

\subsection{Madge}
O \textit{Madge} é uma ferramenta de análise estática de código-fonte utilizada para rastrear, analisar e visualizar grafos de dependência em projetos desenvolvidos em JavaScript e TypeScript. No escopo deste trabalho, o \textit{Madge} é empregado para quantificar o acoplamento estático e o cruzamento de fronteiras (\textit{cross-domain calls}), permitindo auditar o número de importações interdomínio que configuram a conascência de código no monólito base e nas etapas de refatoração.

\subsection{Autocannon}
O \textit{Autocannon} é uma ferramenta de teste de carga e estresse de alta performance voltada para aplicações HTTP em Node.js. Na matriz experimental desta pesquisa, o \textit{Autocannon} atua como o motor de estresse dinâmico para avaliar o desempenho das APIs em cada \textit{snapshot} arquitetural, permitindo mensurar métricas críticas como a latência de resposta e a vazão máxima de requisições por segundo sob condições controladas de concorrência.


